\scanpage{164}
\mspage{83}

est bel et bien un 2-polyèdre -- i.e. s'il satisfait à la condition \deleted{de connexité locale en \unsure{en} chaque sommet}

\[
\begin{aligned}
 &\forall \underset{\substack{\shortparallel\\\{s,t\}}}{a}\in A_0,\ \text{il existe exactement deux }\deleted{\unsure{f\in F_0}}\inserted{u\in S},\\
 &\text{telles que }\quad\{s,u\},\{t,u\}\in A_0,
\end{aligned}\tag*{(19)}
\]

plus la condition de connexité du contour local en tout $s\in S$. S'il en est ainsi, $G_0$ opère bien sûr par autom. de cette structure polyédrale. Si le graphe à homothéties est régulier, il \deleted{\unsure{en ré}} n'en résulte peut-être pas que le 2-polyèdre associé est régulier, \deleted{\unsure{ce}} on dit que le graphe à homothéties est \uline{str.} \uline{régulier} s'il est prépolyédral, et si $G_0$ est transitif (donc simplement transitif) sur l'ens des repères du 2-polyèdre associé.

% slanted
\margin{On a alors $F=\coprod_{i\in\pi}F_i$\ldots\ Dans le cas envisagé\unsure{,| ici,} la structure de multigraphe à hom. est défini par la connaissance de $(S,F\subset\unsure{\mathcal P|\mathfrak P}_3(S),Z\subset\unsure{\mathfrak S|\mathscr S}_S)$.}

Terminologie analogue pour les multigraphes à homothéties \uline{str.} \uline{réguliers} : Alors $G$ est simplement transitif sur l'ens. des repères de $(S,F)$.

J'ai envie de donner une description de ces derniers, en termes de théorie des groupes, en m'appuyant sur la théorie combinatoire des cartes, qui est bien comprise. Je fixe un répère \deleted{$\unsure{(s,t,u)}$ relatif à une structure \unsure{choisie ici} \unsure{constituante}}

\[ (s_0,t_0,u_0)\in\operatorname{Rep}(P_0) \tag*{(20)} \]

pour épingler la situation entièrement. \deleted{\unsure{Du coup}} On veut

% archive sheet number: 163
